%%%%%%%%%%  scan idx 452  |  folio 439  %%%%%%%%%%
En
notant \(E''\) le sous-module de \(E'\) engendré par \(e'_2,\ldots,e'_d\), on voit donc que l'image de \(d:\bigwedge^kE'\longrightarrow\bigwedge^{k-1}E'\) est \(f_1^{(0)}\bigwedge^{k-1}E''\). Par ailleurs, \(f_1^{(0)}\) est régulier dans \(A'\), car sur \(V\) l'idéal inversible \(J'\) est défini par l'équation \(f_1^{(0)}\). Le noyau de \(d:\bigwedge^kE'\longrightarrow\bigwedge^{k-1}E'\) est donc \(\bigwedge^kE''\). Par suite, le \(k\)-ième groupe d'homologie de \(K_\bullet(u)\otimes_AA'\) est \(\bigwedge^kE''/f_1^{(0)}\bigwedge^kE''\); on trouve donc sur \(V\) des isomorphismes:
\begin{equation}\tag{3.2.2}
\underline{\Tor}_i^{\cO_X}(\cO_{X'},\cO_Y)
\xrightarrow{\sim}
(\cO_{X'}/f_1^{(0)}\cO_{X'})\otimes_{\cO_{X'}}
\bigwedge^i(\cO_{X'})^{d-1}
\xrightarrow{\sim}
\bigwedge^i(\cO_{Y'})^{d-1}.
\end{equation}
En particulier, on trouve pour \(i=1\):
\[
\underline{\Tor}_1^{\cO_X}(\cO_{X'},\cO_Y)
\xrightarrow{\sim}
(\cO_{Y'})^{d-1}.
\]
Par suite, on a pour tout \(i\):
\begin{equation}\tag{3.2.3}
\underline{\Tor}_i^{\cO_X}(\cO_{X'},\cO_Y)
\xrightarrow{\sim}
\bigwedge^i\underline{\Tor}_1^{\cO_X}(\cO_{X'},\cO_Y)
\end{equation}

Or d'après 2.2, l'isomorphisme:
\[
\bigwedge(\cO_{Y'})^{d-1}
\xrightarrow{\sim}
\bigoplus_{i\ge0}\underline{\Tor}_i^{\cO_X}(\cO_{X'},\cO_Y)
\]
défini par les isomorphismes (3.2.2) est un isomorphisme d'algèbres. Il en est donc de même pour l'isomorphisme:
\[
\bigwedge\underline{\Tor}_1^{\cO_X}(\cO_{X'},\cO_Y)
\xrightarrow{\sim}
\bigoplus_{i\ge0}\underline{\Tor}_i^{\cO_X}(\cO_{X'},\cO_Y)
\]
défini par les isomorphismes (3.2.3). Pour des raisons d'unicité, c'est alors la restriction à \(V\) de l'homomorphisme (3.2.1), et ce dernier est donc un isomorphisme.

ii) On considère la suite exacte:
\[
0\longrightarrow J\longrightarrow\cO_X\longrightarrow\cO_Y\longrightarrow0,
\]
